\chapter{Podstawy teoretyczne}

Niech $\Sigma$ oznacza alfabet, rozumiany jako skończony zbiór liter. Słowo $w$ jest skończonym ciągiem liter z alfabetu.
Wszystkie słowa nad alfabetem $\Sigma$ oznaczamy jako $\Sigma^*$, a słowo puste jako $\epsilon$. Literą $\Gamma$ oznaczamy alfabet stosowy,
a pusty stos poprzez $\bot$.

\section{Automaty skończone}

\textbf{Deterministycznym automatem skończonym} $A$ dalej DFA $A$ nazywamy krotkę $(\Sigma, Q, q_0, F, \delta)$ w której:
\begin{description}
    \item[$\Sigma$] -- alfabet wejściowy
    \item[$Q$] -- zbiór stanów automatu
    \item[$q_0$] -- stan początkowy
    \item[$F$] -- zbiór stanów akceptujących
    \item[$\delta : Q \times \Sigma \to Q$] -- funkcja przejścia
\end{description}

Domknięciem funkcji $\delta$ nazywamy $\widehat{\delta}: Q \times \Sigma^* \to Q$, zdefiniowaną w następujący sposób:
$$\widehat{\delta}(q, \epsilon) = q$$
$$\widehat{\delta}(q, aw) = \widehat{\delta}(\delta(q, a), w) \text{,\ gdzie\ } a \in \Sigma, w \in \Sigma^*$$
Język $L(A)$ generowany przez automat DFA, to zbiór wszystkich słów $w \in \Sigma^*$, dla których $\widehat{\delta}(q_0, w) \in F$.
Taki język nazywamy językiem regularnym. W dalszej części pracy, gdy z kontekstu będzie jasne, że funkcja przejścia
stosowana jest do słowa, będziemy pisać $\delta$ zamiast $\widehat{\delta}$.

\textbf{Relacja prawostronnej kongruencji.} Relacja $\sim$ jest prawostronną kongruencją, gdy spełnia warunek $\forall_{v_1, v_2 \in \Sigma^*} \forall_{a \in \Sigma}(v_1 \sim v_2 \Rightarrow v_1a \sim v_2a)$

\textbf{Twierdzenie Myhilla-Nerode'a~\cite{Myhill-Nerode}} podaje konieczne i dostateczne warunki, aby język był regularny.
Definiuje ono relację równoważności $\sim_L \subseteq\Sigma^*\times\Sigma^*$, jako $u \sim_L v \Leftrightarrow \forall_{w \in\Sigma^*}(uw \in L \Leftrightarrow vw \in L)$.
Język $L$ jest regularny wtedy i tylko wtedy, gdy relacja $\sim_L$ ma skończenie wiele klas abstrakcji. Liczba stanów minimalnego
automatu DFA rozpoznającego jest równa liczbie klas abstrakcji tej relacji.

\section{Automaty ze stosem}

\textbf{Deterministyczny automat ze stosem} $A$ dalej DPDA $A$ nazywamy krotkę $(\Sigma, Q, q_0, F, \Gamma, \delta)$
w której $\Sigma$, $Q$, $q_0$ i $F$ są zdefiniowane dokładnie jak w przypadku automatu DFA. $\Gamma$ oznacza alfabet stosowy,
a funkcja $\delta : Q\times\Sigma\times\Gamma\to Q\times\Gamma^*$ definiuje przejścia automatu.
Analogicznie do DFA definiujemy domknięcie funkcji $\delta$, jako $\widehat{\delta}: Q\times\Sigma^*\times\Gamma^* \to Q\times\Gamma^*$.
$$\widehat{\delta}(q, \epsilon, s) = (q, s) \text{,\ gdzie\ } s \in \Gamma^*$$
$$\widehat{\delta}(q, aw, sB) = \widehat{\delta}(q', w, ss') \text{,\ gdzie\ } \delta(q, a, B) = (q', s') \text{,\ gdzie\ } B\in \Gamma \text{\ i\ } s, s' \in \Gamma^*$$
Język $L(A)$ generowany przez automat DPDA, to zbiór wszystkich słów $w \in \Sigma^*$, dla których istnieją $q'\in F$ oraz $s \in \Gamma^*$ takie, że $\widehat{\delta}(q_0, w, \bot) = (q', s)$.
Należy zwrócić uwagę, że stos po przeczytaniu słowa $w$ nie musi być pusty, żeby słowo zostało zaakceptowane.


\textbf{Deterministycznym automatem visibly pushdown~\cite{vpl}} $A$ dalej DVPA $A$ nazywamy automat DPDA,
dla którego alfabet wejściowy $\Sigma$ jest podzielony na trzy parami rozłączne zbiory:
$\Sigma_{call}$, $\Sigma_{return}$ i $\Sigma_{local}$,
takie że $\Sigma_{call} \cup \Sigma_{return} \cup \Sigma_{local} = \Sigma$
oraz dla każdego $q \in Q$ zachodzą następujące własności:

$$\forall_{l\in\Sigma_{local}} \forall_{B\in\Gamma} : \delta(q, l, B) = (q', B)$$
Czyli zachowanie automatu po przeczytaniu symbolu $l \in \Sigma_{local}$ nie zależy od symbolu stosowego, ani nie modyfikuje stosu.
Dalej przejścia dla symboli local będziemy oznaczać jako: $\delta(q, l) = q'$ dla uproszczenia notacji.

$$\forall_{c\in\Sigma_{call}} \forall_{B\in\Gamma} : \delta(q, c, B) = (q', BB')$$
Czyli zachowanie automatu po przeczytaniu symbolu $c \in \Sigma_{call}$ również nie zależy od wczytanego symbolu stosowego oraz automat zawsze
wkłada dokładnie jeden dodatkowy symbol na stos. Dalej operacje na symbolach call będziemy oznaczać w~następujący sposób: $\delta(q, c) = (q', B)$.

$$\forall_{r\in\Sigma_{return}} : \delta(q, r, sB) = (q', s)$$
Czyli automat po przeczytaniu symbolu $r \in \Sigma_{return}$ zawsze zdejmuje dokładnie jeden symbol ze stosu. Dalej operacje na symbolach
return będziemy oznaczać w~następujący sposób: $\delta(q, r, B) = q'$

Oznacza to, że klasa automatów DVPA stanowi podklasę DPDA, w której operacja na stosie jest determinowana przez symbol wejściowy.
Model ten zachowuje istotne własności automatów skończonych, takie jak domknięcie ze względu na sumę, przekrój oraz dopełnienie,
przy jednoczesnym umożliwieniu modelowania struktur o zagnieżdżonej naturze.

Dalej litery należące do poszczególnych zbiorów: $\Sigma_{call}$, $\Sigma_{return}$ i $\Sigma_{local}$ będą nazywane odpowiednio
literami typu call, return oraz local.

\textbf{Słowa dobrze zbalansowane.} Słowo $w$ nad alfabetem $\Sigma = \Sigma_{call} \cup \Sigma_{return} \cup \Sigma_{local}$, nazywamy dobrze
zbalansowanym, jeśli w słowie $w$ liczba liter z $\Sigma_{call}$ jest równa liczbie liter z $\Sigma_{return}$ oraz w każdym prefiksie słowa $w$
liczba liter z $\Sigma_{call}$ jest nie mniejsza niż liczba liter z $\Sigma_{return}$.


\textbf{Gramatyką bezkontekstową} $G$ dalej CFG $G$ nazywamy krotkę $(T, N, P, S)$ w której:
\begin{description}
    \item[$T$] -- zbiór symboli terminalnych
    \item[$N$] -- zbiór symboli nieterminalnych
    \item[$P$] -- zbiór reguł typu $L \rightarrow R$, gdzie $L \in N$ oraz $R \in (T \cup N)^*$
    \item[$S$] -- $S \in N$ jest elementem początkowym
\end{description}
Słowo $w \in T^*$ jest generowane przez gramatykę, jeśli można je otrzymać z symbolu początkowego $S$ przez skończony
ciąg produkcji z $P$. Zbiór wszystkich słów generowanych przez gramatykę $G$ oznaczamy jako $L(G)$.

\section{Systemy przepisywania słów}

\textbf{System Przepisywania Słów (ang. String Rewriting System)} $R$ dalej SRS $R$ nazywamy zbiór reguł typu $l \to r$,
gdzie $l$ oraz $r$ są słowami nad pewnym alfabetem $\Sigma$. Jeden krok przepisywania względem systemu $R$ oznaczamy przez $\to_R$.
Dla reguły $l \to r \in R$ oraz słów $x,y \in \Sigma^*$ zachodzi $xly \to_R xry$.
Przez $\to_R^*$ oznaczamy zwrotno-przechodnie domknięcie relacji $\to_R$.
W tej pracy systemy SRS będą wykorzystywane jako więzy opisujące zależności między słowami względem rozważanego języka.

\textbf{Spójność systemu SRS z językiem.} System SRS $R$ jest spójny z językiem $L$ wtedy i tylko wtedy, gdy dla wszystkich słów $s, t \in \Sigma^*$,
jeśli $s \to^*_R t$, to $s \in L \Leftrightarrow t \in L$.

\textbf{Reguły ze zmiennymi.} Niech $\mathcal{X}$ będzie skończonym zbiorem zmiennych. Regułą przepisywania słów ze zmiennymi
nazywamy regułę $\alpha \to \beta$, gdzie $\alpha, \beta \in (\Sigma \cup \mathcal{X})^*$ oraz:
\begin{enumerate}
    \item Każda zmienna $X \in \mathcal{X}$ występuje co najwyżej raz w~$\alpha$ oraz co najwyżej raz w~$\beta$.
    \item Każda zmienna $X \in \mathcal{X}$, która występuje w $\beta$, występuje również w $\alpha$.
\end{enumerate}
Dla każdej zmiennej we wszystkich jej wystąpieniach podstawiane jest to samo słowo. System SRS zawierający reguły ze zmiennymi
nazywamy systemem SRS ze zmiennymi.

\textbf{Instancją reguły ze zmiennymi} nazywamy zwykłą regułę SRS otrzymaną przez podstawienie słów dobrze zbalansowanych za zmienne
w regule ze zmiennymi.

\textbf{Reguły zachowujące dopasowanie.} W przypadku języków visibly pushdown regułę ze zmiennymi $\alpha \to \beta$ nazywamy
zachowującą dopasowanie, jeśli każda jej instancja jest regułą zachowującą dopasowanie.
Dla zwykłej reguły $l \to r$ oznacza to, że dla wszystkich słów $w,w' \in \Sigma^*$ takich, że $w=xly$ oraz $w'=xry$,
relacja dopasowania pomiędzy symbolami typu call występującymi w $x$ oraz symbolami typu return występującymi w $y$ jest
taka sama dla słów $w$ oraz $w'$. System SRS nazywamy zachowującym dopasowanie, jeśli wszystkie należące do niego reguły
zachowują dopasowanie.

\textbf{Reguły dobrze zbalansowane.} Szczególnym przypadkiem reguł zachowujących dopasowanie są reguły dobrze
zbalansowane. Regułę $\alpha \to \beta$ nazywamy dobrze zbalansowaną, jeśli zarówno słowo $\alpha$, jak i słowo $\beta$ po usunięciu
zmiennych są dobrze zbalansowane. System SRS składający się wyłącznie z reguł dobrze zbalansowanych nazywamy dobrze zbalansowanym.